\subsection{Straight-wall sectors and their dual networks}
\label{sector:radial-setup}

Fix an integer $k\geq1$.  Take $k$ consecutive copies of a coordinate
quadrant, glue their intervening radial sides in order, and do not identify
the first and last sides.  All the constructions in this subsection take
place on this noncyclic angular cover.  For integers $1\leq r<R$, let
$H_k(r,R)$ be the mesh-one graph between the sup-norm levels $r$ and $R$,
including all vertices and edges in the closed region.  Its two level
arcs are denoted by $S_r$ and $S_R$, and its two remaining boundary arcs
are its \emph{flanks}.  The underlying surface is a closed disk.  For
$k=4$ the two banks of the coordinate slit remain distinct; the same
convention applies when $k>4$.  Thus local square-grid coordinates are
used on an abstract disk, not on a cyclically identified annulus.
Set
\begin{equation}\label{sector:Q-definition}
 Q_k(r,R)=g_{H_k(r,R)}(S_r,S_R),
 \qquad
 \mathcal H(p)=\frac{1-p}{1-up}.
\end{equation}
The two level arcs in this definition are wired separately.  Constants
below may depend on the fixed $k$ and on $q$, but not on the radii.

\begin{proposition}[Radial network decay]\label{prop:sector-network}
For every fixed finite $k\geq1$, every lattice rotation or reflection of
the construction, and every translate of its center to a lattice site,
\begin{equation}\label{sector:network-decay}
 \lim_{L\to\infty}\ \sup_{\substack{1\leq r<R\\R/r\geq L}}
 Q_k(r,R)=0.
\end{equation}
This is a separately wired two-terminal network statement on the
noncyclic disk sector.
\end{proposition}

We first record the geometry and duality needed for its proof.  Let
$C_k(r,R)$ be the cell-center graph of the sector: its vertices are the
centers of its square cells, and neighbors are joined across interior
edges.  To this graph add two distinct hubs, one for each flank, with an
odds-$v$ spoke dual to every edge on that flank.  Denote the connection
probability of the two hubs by $P_k^*(r,R)$.

\begin{lemma}[Exact split-hub duality]\label{sector:duality}
One has
\begin{equation}\label{sector:duality-equation}
 P_k^*(r,R)=\mathcal H(Q_k(r,R)),\qquad
 \mathcal H(\mathcal H(p))=p.
\end{equation}
The same assertion holds for a cell-center passage between two
concentric layer cuts, using its own cell dual and its two flank hubs.
\end{lemma}

\begin{proof}
Contract each level arc separately.  Edges along a contracted arc are
loops and contribute factors independent of the remaining configuration,
so they may be omitted.  The ordinary plane dual has precisely the
interior cell graph just described, with a \emph{single} exterior hub
carrying the spokes from both flanks.  Its odds are $q/v=v$.
Euler's formula gives the FK dual law with these odds.  The disk
path--cut alternative identifies absence of the primal level-to-level
connection with an open dual traversal from one flank to the other
through the cells, without using the common hub as a shortcut.

Split the common hub into its two flank hubs and let $p$ be their
connection probability in the split law.  Joining the hubs changes the
component count by one on separation and by zero on connection.
Consequently the joined law gives the traversal event probability
\[
 \frac{p}{p+t(1-p)}=\frac{q p}{1-up}.
\]
This is $1-Q_k(r,R)$.  Solving for either probability proves
\eqref{sector:duality-equation}; direct substitution proves that
$\mathcal H$ is involutive.  This reasoning uses only the disk boundary
order and the square-cell incidence structure, so applies to the
cell-center passage as well.  No Euclidean identification across a slit
or between different sheets of the angular cover is made.
\end{proof}

The next lemma turns rectangle crossings into a connection along a
boundary route.  It will be used twice: along the outer rim to bound the
split-dual test, and around an end gap to extend a measured actual arm.
The latter use requires the reflex-corner case as well as convex corners.

\begin{lemma}[Boundary chords along a rectilinear route]
\label{sector:blank-boundary-chain}
Consider a square-grid disk graph, in ordinary coordinates or on a
noncyclic angular cover, with two accessible boundary intervals $I$ and
$J$, separately wired for the test and with no other ties.  Suppose that
a boundary route from $I$ to $J$
has a bounded number of straight portions of comparable macroscopic
length, positive inward clearance of comparable size along each
portion, and right-angle convex or reflex corners.  The fitting inward
rectangular neighborhoods along straight portions and inward square
neighborhoods at convex corners must contain every grid vertex and
nearest-neighbor edge of their closed regions, including boundary edges,
with ordinary odds $v=\sqrt q$.  At a reflex corner, require the same
complete ordinary grid on a fitting closed square of comparable size
with its exterior quadrant removed.  Suppose also that $I$ and $J$
have positive length on this same scale.  Then
\[
 g(I,J)\ge c>0,
\]
where $c$ depends only on $q$, the bound on the number of straight
portions, and the fixed length and clearance ratios.
The estimate remains valid after a cofacial extension for
which every added edge is accessible from a test-edge face by
crossing only added edges.
\end{lemma}

\begin{proof}
All the comparisons here use auxiliary interval chains in narrow
exterior caps and an exterior terminal test edge.  When a rectangle
is flush to a continuous exposed boundary stretch, every unwanted
edge can be reached from a test-edge face across only unwanted edges.
Successive face expansion therefore peels the unwanted edges, and
cofacial Rayleigh gives the rectangle-to-disk comparison in the
required direction.  Interval caps can be chosen with no physical
edge in their thin pockets, so they do not obstruct the peeling.

Proposition~\ref{prop:rectangle-rsw} first gives chords from a fixed
positive-length patch to the opposite side of a fitting subrectangle.
For two patches on adjacent sides, the two patch-to-opposite-side
chords interlace.  Lemma~\ref{lem:interlacing} therefore gives a
positive lower bound for their connection.  For two patches on the
same side, first connect each toward the opposite side.  Interlace
each of these chords with a transverse rectangle crossing, choosing
the right-hand endpoint for the left patch and the left-hand endpoint
for the right patch.  The resulting two chords again interlace and
connect the original patches.  When two input full-side electrodes meet
at a corner, use
Corollary~\ref{rect:interlacing-abutments}; the two patches to be
connected are disjoint and stay away from that corner.  All rectangles
have fixed aspect ratios and positive relative patch sizes.  At a convex
corner the same adjacent-side construction takes place in an inward square.

Here is the corresponding construction at a reflex corner.  In local
coordinates take $[-h,h]^2$ with its southwest quadrant removed.
Let the two boundary patches lie on its negative horizontal and
negative vertical rays, away from the corner.  In the upper rectangle,
connect the first patch to an interval in the upper portion of the
right wall.  In the right rectangle, connect the second patch to an
interval in the right portion of the top wall.  The adjacent-side
construction just given bounds both probabilities below.  Extend
each rectangle cofacially to the clipped square.  The four intervals
occur in the cyclic order
\[
 \text{negative horizontal ray},\quad
 \text{negative vertical ray},\quad
 \text{right wall},\quad \text{top wall}.
\]
The two chords consequently interlace.  A further application of
Lemma~\ref{lem:interlacing} connects the two original patches with
a fixed positive lower bound.

Subdivide the boundary route into a fixed finite number of sufficiently
short comparable pieces, putting disjoint ordered patches
$I_0,\ldots,I_N$ on their middle portions, with $I_0\subset I$
and $I_N\subset J$.  Choose the subdivision so a two-piece span
meets at most one corner and fits one of the preceding local
constructions.  Those constructions give
\begin{equation}\label{sector:skip-chords}
 g(I_j,I_{j+2})\ge c_0>0.
\end{equation}
Interlace the chord from $I_0$ to $I_j$
with the chord from $I_{j-1}$ to $I_{j+1}$, beginning with $j=2$.
After a bounded number of applications this gives
\begin{equation}\label{sector:route-chord}
 g(I_0,I_N)\ge c_1>0.
\end{equation}
Enlarge the terminal intervals to $I,J$ by
adjacency.  Finally, the face-expansion comparison in the first
paragraph proves the stated extension property.
\end{proof}

\begin{lemma}[Fixed-geometry bounds for the two dual tests]
\label{sector:fixed-geometry}
There exist $c_0,c_1>0$ and finite $L_0,m_0$ such that
\begin{enumerate}[label=\textup{(\roman*)}]
\item $P_k^*(r,R)\geq c_0$ whenever $R/r\geq L_0$;
\item in the spoke-free cell graph consisting of centers with
\[
 m+\tfrac12\leq |z|_\infty\leq2m-\tfrac12
\]
on the same $k$-quadrant cover, the network connection between its two
whole layer cuts, separately wired, is at most $1-c_1$ for $m\geq m_0$.
\end{enumerate}
\end{lemma}

\begin{proof}
For (i), use the convex-corner case of
Lemma~\ref{sector:blank-boundary-chain} along the outer polygon rim of
$C_k(r,R)$,
starting and ending on flank patches at distances comparable to $R$
from the center.  When $R/r$ is sufficiently large, the route stays a
fixed fraction of $R$ away from the inner rim.  On each sheet its
straight pieces have lengths comparable to $R$, and the local inward
rectangles and corner squares have sides a small fixed fraction of
$R$.  Crossing an intervening quadrant seam does not change the local
square-grid geometry.  There are only finitely many pieces for fixed
$k$.  The two banks are always treated as separate disk sides.
The lemma therefore gives a deterministic-contact network lower bound
between the two flank patches.

Replacing such a deterministic contact by its stochastic spokes loses
at most the factor
\[
 \eta=\frac{v}{1+v}>0.
\]
Indeed, with the other tested terminal fixed, the rooted formula of
Lemma~\ref{lem:rooted} writes the new isolated hub's connection
probability as
\[
 \EE_\mu\bigl[1-(1+v)^{-N}\bigr]
 \geq \eta\,\mu(N>0),
\]
where $N$ counts rooted-multiset multiplicity on the contact patch.
The right-hand hitting probability is the deterministic-contact
network.  Perform this comparison for the two hubs in turn.  Adding
their other spokes does not decrease their connection probability by
adjacent Rayleigh.  This proves (i), with a factor $\eta^2$.

For (ii), contract the two circular layer cuts of the stated
spoke-free passage and apply Lemma~\ref{sector:duality}.  The split
dual has ordinary square-grid interiors and stochastic flank spokes.
Its outer boundary route has the same local form as the one just
used.  Taking a cell dual shifts its rows and endpoints by bounded
mesh distances; straight lengths and available inward depths remain
comparable to $m$, and successive corner rows still meet at right
angles.  For $m$ sufficiently large the fixed clearances used above
absorb all these shifts.  The split dual flank-to-flank probability
is therefore at least a positive constant $c_4$.  Its complementary
radial probability is at most $\mathcal H(c_4)<1$ by exact duality.
Take $c_1=1-\mathcal H(c_4)$.
\end{proof}

\subsection{A noncrossing partition obstruction}
\label{sector:partition-obstruction}

The obstruction comes from a concrete finite graph inside a long
sector.  To see the problem it must exclude, suppose the extremal
radial probabilities along some sequence tend to a number
$M\in(0,1)$, and put $h=\mathcal H(M)$.  Choose a fixed number of
radial bands whose ratios, intervening gaps, and distances from the
two ends all tend to infinity.  On each flank of the interior cell
graph, attach a separate hub through the ordinary spokes of each
band.  Write $a_i$ and $b_i$ for the two hubs of band $i$, in increasing
radial order.  Their open connection partition is noncrossing in
the boundary order
\[
 a_1,\ldots,a_n,b_n,\ldots,b_1.
\]

Two comparisons will constrain this finite partition.  If a group of
first-flank hubs and a group of second-flank hubs contain a matched
pair $a_i,b_i$, their connection after the two groups are separately
merged is bounded below by the isolated split-dual test of band $i$
and above by the full split-dual test.  Both bounds tend to $h$.
By contrast, two ordered groups on one flank can be enlarged to the
full layer cuts of an unused intervening gap, whose network probability
is at most $1-c_1$.  Terminal enlargement also bounds this group test
below by the connection probability of any one hub from each group.
The underlying hub law also retains the pressure
inequality of Theorem~\ref{thm:pressure}.  The completion of the proof
will verify these comparisons from the permitted finite edge changes.
The next lemma rules out the infinite paired-list law that their
limiting consequences would produce by finite partition compactness.

For a partition $\pi$ and a finite vertex set $D$, let $k_D$ count the
blocks meeting $D$.  For disjoint nonempty finite sets $D,F$, put
\[
 C_\pi(D,F)=\one\{\text{a block of $\pi$ meets both $D$ and $F$}\},
 \qquad r_\pi(D,F)=k_{D\cup F}-2+C_\pi(D,F).
\]
The integer $r_\pi(D,F)$ is precisely the rank gain on separately
merging $D$ and $F$: after these mergers the blocks meeting their
union number $2-C_\pi(D,F)$.  Define
\begin{equation}\label{sector:merger-functional}
 G(D,F)=\frac{\EE[t^{r_\pi(D,F)}C_\pi(D,F)]}
                   {\EE[t^{r_\pi(D,F)}]},
 \qquad K(D)=\log\EE[t^{k_D}].
\end{equation}
These definitions depend only on the restriction of $\pi$ to
$D\cup F$, even if the underlying partition is infinite.
For an FK connection partition, $G(D,F)$ is exactly the network
connection probability after separately merging the two groups:
$t^{r_\pi(D,F)}$ is the change of component-counting weight.
Thus $G$ records the group tests in the finite fan graph just
described.  The lemma makes no initial exchangeability assumption.

\begin{lemma}[Paired-list obstruction]\label{lem:sector:paired-lists}
There is no random partition of
$\{a_i,b_i:i\in\N\}$ with all the following properties, for fixed
$t>1$, $h\in(0,1)$, and $c>0$.
\begin{enumerate}[label=\textup{(\roman*)}]
\item Each restriction to the first $n$ pairs is noncrossing in the
cyclic order $a_1,\ldots,a_n,b_n,\ldots,b_1$.
\item For disjoint finite sets $D,F$,
\begin{equation}\label{sector:partition-pressure}
 K(D)+K(F)-K(D\cup F)\geq0.
\end{equation}
\item For nonempty finite index sets $U,V$ with $U\cap V\ne\varnothing$,
\begin{equation}\label{sector:matched-tests}
 G(a_U,b_V)=h,
 \qquad a_U=\{a_i:i\in U\},\quad b_V=\{b_i:i\in V\}.
\end{equation}
\item If $U,V$ are nonempty and $\max U<\min V$, then for every
$i\in U$ and $j\in V$,
\begin{equation}\label{sector:same-side-tests}
 \PP(a_i\sim a_j)\leq G(a_U,a_V)\leq1-c.
\end{equation}
\end{enumerate}
\end{lemma}

\begin{proof}
\emph{How the obstruction works.}
We first replace the partition law by a spreadable one and express it
as an iid mixture of labeled pairs.  Noncrossing leaves only private
matched bonds and at most one repeated label on either side.  The
merger identities and pressure inequality leave two pure cases; in
the private case the same pressure inequality makes the bond probability
constant.  The matched tests exclude that case, and the same-side
tests exclude the repeated-label case.

\emph{Making the law spreadable.}
For each $m$ there are finitely many partitions of $2m$ labeled
vertices.  Color an increasing $m$-tuple of pair indices by a finite
quantization of its vector of restriction probabilities.  Iterated
infinite Ramsey \cite[Theorem~A]{Ramsey1930} gives nested infinite sets
$J_j$ such that, for every
$m\leq j$, these vectors differ by at most $2^{-j}$ between any two
increasing $m$-tuples in $J_j$.  Restrict the law to the first $j$
indices of $J_j$ and take a projective subsequential limit, using
compactness of each finite partition simplex.  The resulting law is
invariant under increasing injections of the pair indices.  Every
condition in the statement survives: it involves finitely many
partition probabilities, the merger quotients have positive
denominators, and increasing restrictions preserve the specified
boundary order and the separation of $U,V$.  We may therefore assume
this spreadability from now on.

Give each partition block an independent uniform label in $[0,1]$.
There are countably many blocks, and their labels are almost surely
distinct.  Let $X_i,Y_i$ label $a_i,b_i$, and put $Z_i=(X_i,Y_i)$.
This sequence is spreadable on the compact space $[0,1]^2$.
The spreadable form of the de Finetti representation is due to
Ryll--Nardzewski \cite{RyllNardzewski1957}; we include the argument for
this compact label space.  For a bounded continuous real
function $f$ and two disjoint index blocks $B,C$, each of length $L$,
spreadability gives
\begin{equation}\label{sector:block-average}
 \EE\left[\left(\frac1L\sum_{i\in B}f(Z_i)
                    -\frac1L\sum_{i\in C}f(Z_i)\right)^2\right]
 =\frac2L\bigl(\EE f(Z_1)^2-\EE[f(Z_1)f(Z_2)]\bigr).
\end{equation}
For functions $f_1,\ldots,f_m$, average their product using one index
in each of $m$ successive length-$L$ blocks.  Its expectation equals
$\EE\prod_j f_j(Z_j)$ by spreadability.  Replacing every block average
by the average of the same function over the first block changes this
expectation by $O(L^{-1/2})$, by \eqref{sector:block-average} and a
telescoping product estimate.  The laws of the empirical measures of
the first block have a subsequential weak limit on the compact space
of probability measures on $[0,1]^2$.  Calling the resulting random
measure $\nu$, we obtain
\[
 \EE\prod_{j=1}^m f_j(Z_j)=\EE\prod_{j=1}^m\nu(f_j).
\]
Products of continuous functions determine the finite-dimensional
laws.  We can consequently realize the labeled pairs as iid with
common law $\nu$, conditionally on $\nu$.

\emph{Noncrossing restricts the directing law.}
Almost every directing law gives a noncrossing partition with
conditional probability one, simultaneously for all finite
restrictions.  Its first marginal has at most one atom: if two
distinct labels $x,y$ had positive mass, the event
$X_1=x,X_2=y,X_3=x,X_4=y$ would have positive probability and would
cross.  The second marginal has the same property.  Denote their
possible atom masses by $\alpha,\beta$, putting the mass equal to zero
for a diffuse marginal.

A diagonal tie $X_i=Y_i$ at a label other than a common marginal atom
has positive probability only if both marginals are diffuse.  For
example, if the first marginal has atom $x$ of mass $\alpha>0$ and
such diagonal ties have mass $\delta>0$, three independent pairs
produce $X_1=X_3=x$ and $X_2=Y_2\ne x$ with probability
$\alpha^2\delta>0$.  The two blocks interlace in the disk boundary
order, which is impossible.  The second-side argument is identical.
Let
\[
 D_0=\{\alpha=\beta=0\},\qquad d_0=\PP(D_0).
\]
On $D_0$, labels from different pairs are almost surely distinct,
since each marginal is diffuse and different pairs are independent.
Thus the only ties are private matched bonds, conditionally
independent with a probability $\gamma=\nu\{(x,x):x\in[0,1]\}$.
Outside $D_0$ every repetition uses a marginal atom, and no diffuse
matched bonds occur.  The two atoms, if both exist, need not be equal.

\emph{The merger weights.}
Write
\begin{equation}\label{sector:A-definition}
 w(m)=t^{-m+\one\{m>0\}},\qquad
 A_j(x)=t(1-ux)^j+(1-t)(1-x)^j,
 \qquad u=1-t^{-1}.
\end{equation}
If $M$ is binomial with parameters $j,x$, then
\begin{equation}\label{sector:A-bounds}
 \EE w(M)=A_j(x),\qquad A_1(x)=1,\qquad
 (1-ux)^j\leq A_j(x)\leq t(1-ux)^j,\qquad A_j(x)\leq1.
\end{equation}
The first identity follows by separating $M=0$ and summing
$t\EE[t^{-M}]$ on $M>0$; the bounds also follow directly from the
formula, or from $0<w(m)\leq1$.

Test $k$ first-side labels and $s$ second-side labels with exactly
$l$ overlapping pair indices, where $1\leq l\leq k$ is fixed and
$s\to\infty$.  Normalize the reweighting by $t^{k+s-2}$.
On $D_0$, if $m$ of the $l$ matched bonds occur, the union has
$k+s-m$ blocks and $C=\one\{m>0\}$, so the normalized weight is
$w(m)$.  The conditional denominator and numerator are respectively
\begin{equation}\label{sector:private-weights}
 A_l(\gamma),\qquad A_l(\gamma)-(1-\gamma)^l.
\end{equation}
On $\{\beta=0<\alpha\}$ there are no cross-group ties, and the
denominator is $A_k(\alpha)$ while the numerator is zero.
On $\{\beta>0\}$ the normalized weight tends to zero in conditional
expectation: impose the second-group merge first.  Its normalized
factor is $w(M_s)$, where $M_s$ counts occurrences of its atom; the
remaining first-group merge has rank gain at most $k-1$.
Thus the whole normalized weight is at most $w(M_s)\leq1$, and
$M_s\to\infty$ almost surely.  Bounded convergence applies also
after averaging over $\nu$.

Equation~\eqref{sector:matched-tests}, multiplied by its denominator,
therefore yields, for every $k\geq l$,
\begin{equation}\label{sector:half-atomic-equation}
 \EE\bigl[\one_{D_0}\{A_l(\gamma)-(1-\gamma)^l\}\bigr]
 =h\left(\EE[\one_{D_0}A_l(\gamma)]
          +\EE[\one_{\{\beta=0<\alpha\}}A_k(\alpha)]\right).
\end{equation}
Since $A_k(\alpha)\to0$ on $\alpha>0$, take $l=1$, let $k\to\infty$,
and compare with $k=1$.  As $A_1=1$ and $h>0$, this proves
$\PP(\beta=0<\alpha)=0$.  Reversing the two sides proves
$\PP(\alpha=0<\beta)=0$.  No division by a limiting denominator has
been used.  Hence outside $D_0$ both atom masses are positive.

\emph{Eliminating a mixed diffuse/atomic law.}
Choose $D=\{a_1,\ldots,a_s\}$ and
$F=\{b_{s+1},\ldots,b_{2s}\}$.  On $D_0$ all selected labels are
distinct.  Outside $D_0$ the number of repetitions of each marginal
atom tends to infinity.  Consequently each of
\[
 \EE[t^{k_D-s}],\qquad \EE[t^{k_F-s}],\qquad
 \EE[t^{k_{D\cup F}-2s}]
\]
tends to $d_0$.  For the last term, use
$k_{D\cup F}\leq k_D+k_F$; all three integrands are bounded by one.
If $0<d_0<1$, this gives
\[
 K(D)+K(F)-K(D\cup F)\longrightarrow\log d_0<0,
\]
contrary to \eqref{sector:partition-pressure}.

\emph{The purely diffuse case.}
If $d_0=1$, put $V=1-u\gamma$ and use the two disjoint vertex sets
$D=\{a_1,b_1\}$ and $F=\{a_2,b_2\}$.  Conditional independence of
the private bonds gives
\[
 \EE t^{k_D}=t^2\EE V,\qquad
 \EE t^{k_{D\cup F}}=t^4\EE V^2.
\]
Pressure positivity forces $(\EE V)^2\geq\EE V^2$, so $\gamma$ is
constant.  The singleton matched test in
\eqref{sector:matched-tests} gives $\gamma=h>0$.  For groups of
size $s$ with all indices paired, \eqref{sector:private-weights}
gives failure probability
\[
 \frac{(1-\gamma)^s}{A_s(\gamma)}
 \leq\left(\frac{1-\gamma}{1-u\gamma}\right)^s\longrightarrow0.
\]
Their connection probabilities cannot all equal $h<1$.

\emph{The positive-atom case.}
If $d_0=0$, the first marginal has exactly one atom of mass
$\alpha>0$ almost surely.  Two distinct first-side vertices connect
with probability $\EE\alpha^2>0$.  For ordered-separated groups of
sizes $s,k$, their atom counts are independent conditional on $\nu$.
There is at most one block shared by the groups.  Their normalized
merger weight is the product of their two single-group factors, and
connection requires a positive atom count in both.  Thus
\begin{equation}\label{sector:atomic-group-ratio}
 G(a_U,a_V)=
 \frac{\EE[(A_s(\alpha)-(1-\alpha)^s)
                (A_k(\alpha)-(1-\alpha)^k)]}
      {\EE[A_s(\alpha)A_k(\alpha)]}.
\end{equation}
By \eqref{sector:same-side-tests} this lies in
$[\EE\alpha^2,1-c]$ for every $s,k$.

If $\operatorname{ess\,inf}\alpha=0$, fix $k=1$.
The probability laws weighted by $A_s(\alpha)$ concentrate at zero.
Indeed, for any $0<\epsilon<1$, the positive-mass set
$\{\alpha\leq\epsilon/2\}$ contributes at least
$\PP(\alpha\leq\epsilon/2)(1-u\epsilon/2)^s$ to the denominator,
whereas $\{\alpha\geq\epsilon\}$ contributes at most
$t(1-u\epsilon)^s$.  Their ratio tends to zero.
Since $A_1=1$ and $A_1-(1-\alpha)=\alpha$, the ratio in
\eqref{sector:atomic-group-ratio} is at most the expectation of
$\alpha$ under this weighted law and therefore tends to zero.
This contradicts its positive lower bound.
If instead $\alpha\geq a>0$ almost surely, then uniformly
\[
 \frac{(1-\alpha)^s}{A_s(\alpha)}
 \leq\left(\frac{1-a}{1-ua}\right)^s\longrightarrow0.
\]
Taking $k=s$ in \eqref{sector:atomic-group-ratio} makes the ratio tend
to one, contradicting its strict upper bound.  All possibilities for
$d_0$ have been excluded.
\end{proof}

\subsection{Separated radial fans and completion of network decay}
\label{sector:fan-comparison}

\begin{proof}[Proof of Proposition~\ref{prop:sector-network}]
For the fixed $k$ and orientation, the nonincreasing supremum in
\eqref{sector:network-decay} has a limit $M\geq0$.
Lemma~\ref{sector:fixed-geometry}(i) and exact duality imply
$M\leq\mathcal H(c_0)<1$.  Suppose $M>0$.  Choose sectors with
$R/r\to\infty$ and $Q_k(r,R)\to M$, and put
\[
 h=\mathcal H(M)\in(0,1).
\]
Fix $n$.  Inside each sector choose integer levels satisfying
\begin{equation}\label{sector:separated-levels}
 r\ll s_1^-\ll s_1^+\ll s_2^-\ll\cdots
       \ll s_n^-\ll s_n^+\ll R,
\end{equation}
where every consecutive ratio tends to infinity along the chosen
sequence.  For example, divide the logarithmic interval from $r$ to
$R$ into $2n+1$ equal parts and round the resulting levels; rounding
does not affect these ratios.  In particular the interior levels
tend to infinity.

Use the whole interior cell graph $C_k(r,R)$ as a base.  On the first
flank put distinct hubs $a_1,\ldots,a_n$, with $a_i$ having precisely
the odds-$v$ spokes dual to flank edges between $s_i^-$ and $s_i^+$.
Put analogous hubs $b_1,\ldots,b_n$ on the second flank.  There are
no other spokes.  Draw each fan in a narrow exterior strip along its
own boundary interval.  These strips are disjoint, so the network
partition of the hubs is noncrossing in the order
$a_1,\ldots,a_n,b_n,\ldots,b_1$.

For a matched pair $a_i,b_i$, retain just the cell subband between
$s_i^-$ and $s_i^+$ and its two fans.  Its connection probability is
exactly $\mathcal H(Q_k(s_i^-,s_i^+))$.  Extending it to the full
interior cell graph, and then adding the remaining separate fans,
does not decrease this tested network probability.  Here is the
cofacial comparison.  Draw the test edge joining the two hubs in the
exterior of the disk.  The complement of the closed subband has an
inner and an outer portion.  Each unwanted interior-grid edge has
access through its portion to the corresponding inner or outer rim,
and thence to a test-edge face, crossing only unwanted edges.  The
two retained fan strips occupy only their stated flank intervals;
their pockets contain no unwanted interior edges and they do not
close these complementary openings.  Face expansion therefore
deletes all unwanted edges in an order cofacial with the test edge.
Cofacial Rayleigh gives the claimed extension comparison.  Each
other fan lies on a disjoint boundary stretch; peel its edges
successively from the main exterior face and reverse this order to
add the fan.  The same cofacial comparison applies at each addition.
All these are comparisons of finite-odds edges, so no unproved
boundary-condition monotonicity is involved.  Thus
\begin{equation}\label{sector:matched-lower}
 \liminf\PP(a_i\sim b_i)\geq h,
\end{equation}
because each subband ratio tends to infinity and the defining
supremum for $M$ bounds its $Q_k$ from above by $M+o(1)$.

Conversely, merge all first-flank hubs into one terminal and all
second-flank hubs into another, then add every missing flank spoke.
Each merger can be made by edges incident to the corresponding
tested terminal, and the added spokes are incident there too.
Adjacent Rayleigh shows that the network probability does not decrease
at any step.  The result
is the full split dual network, with connection probability
$\mathcal H(Q_k(r,R))\to h$.
More generally, for any nonempty $U,V\subset\{1,\ldots,n\}$ with
$U\cap V\ne\varnothing$, choose $i\in U\cap V$.  Enlarge the
matched terminals $a_i,b_i$ to $a_U,b_V$, and then to both full flank
terminals.  The two preceding comparisons squeeze the merger test
to give
\begin{equation}\label{sector:finite-overlap-limit}
 G(a_U,b_V)\longrightarrow h.
\end{equation}
Here $G$ is exactly \eqref{sector:merger-functional}: the FK change
of component count on identification is its rank-gain reweighting.

There is also a uniform strict upper bound for ordered-separated
groups on the first flank.  Suppose $\max U<\min V$.  In a gap
between consecutive band intervals lying between these groups,
choose an integer $m\to\infty$ so that the two cell layers
$m+\tfrac12$ and $2m-\tfrac12$ and the entire passage between them
lie strictly inside the gap.  This is possible by
\eqref{sector:separated-levels}.  Enlarge the inner tested terminal
to the whole inner layer and the outer terminal to the whole outer
layer.  All vertices, hubs, and fans below the gap may be absorbed
into the first terminal, and all those above it into the second.
No fan crosses the gap.  Every enlargement is by links incident to
a tested terminal and does not decrease its connection probability.
After these identifications the two exterior portions attach to
the passage only at its respective terminal; their partition
functions factor out.  The remaining network is precisely the pure
spoke-free cell-layer passage of
Lemma~\ref{sector:fixed-geometry}(ii).  Therefore
\begin{equation}\label{sector:finite-same-side}
 \PP(a_i\sim a_j)\leq G(a_U,a_V)\leq1-c_1
 \quad(i\in U,\ j\in V)
\end{equation}
once the original sector is sufficiently far along the sequence.
The lower bound is the direct terminal-enlargement comparison in
the original finite hub graph.

Finally Theorem~\ref{thm:pressure}, applied to that finite graph,
gives \eqref{sector:partition-pressure} for every pair of disjoint
hub sets.  For fixed $n$, pass to a subsequential limit of its finite
partition law.  Equations~\eqref{sector:finite-overlap-limit} and
\eqref{sector:finite-same-side}, noncrossing, and pressure positivity
survive.  Let $n$ increase and take another projective subsequential
limit on the two infinite lists.  This is legitimate because each
finite partition space is finite, and all constraints on a fixed
restriction hold for all sufficiently large $n$.  The resulting law
satisfies every hypothesis of Lemma~\ref{lem:sector:paired-lists},
with $c=c_1$, a contradiction.  Hence $M=0$.

For each fixed larger $k$ the outer route has more, but still
boundedly many, straight pieces and convex corners.  The local
comparisons, cell duality, radial band order, and gap contraction
are unchanged on the noncyclic angular cover.  Thus the same proof
establishes the claimed range of $k$, with constants allowed to
depend on it.  Rotations, reflections, and lattice translations
preserve all graph incidences and odds.
\end{proof}
